\chapter{警报信号}
% 完整版：chapters/14_pain.tex

\pencilfig{14}{\pencilart{14}}{疼痛是火灾警报器，而不是温度计。}

其他所有感觉报告的都是世界是什么样的。疼痛不是。它报告的是有东西在威胁身体，而它的音量在很大程度上由机器自己设定，而不是由刺激决定。疼痛是火灾警报器，而不是温度计。

痛觉传感器只对强烈的刺激起反应：高于四十多度的热、强大的压力、腐蚀性物质。它们的适应比触觉传感器弱得多——只要威胁还在，警报就不该停。

\section*{两根导线}

回想你撞到手指的时候。先是尖锐、精确的疼：在哪里，什么时候。片刻之后，是钝的、弥漫的、持久的疼：有多糟糕。这是两根不同的导线。一根快，它的信号几十毫秒就能送到。另一根慢，它的信号要走大约一秒。

\section*{闸门}

为什么撞到以后我们会去揉？在脊髓里，疼痛信号要经过一道“闸门”。触摸会唤醒一个抑制性神经元，把闸门掩上。所以揉一揉确实管用——但只对中等程度的疼痛有效；剧烈的疼痛照样通过。至于剧痛会同时把这个抑制性神经元关掉，这是最初的闸门理论的一部分，目前还没有直接证据。

\pencilfig{126}{%
% gate: the pain wire drives the relay neuron; touch wakes an inhibitory neuron that closes the gate
\draw[pen=pcAmber] (0,1.35) -- (3.05,1.0); \draw[arr=pcAmber] (2.8,1.03) -- (3.08,1.0);
\node[lbl, text=pcAmber!70!black, anchor=east] at (-0.05,1.35) {疼痛};
\draw[pen=pcBlue] (0,-0.15) -- (1.75,-0.15); \draw[arr=pcBlue] (1.5,-0.15) -- (1.78,-0.15);
\node[lbl, text=pcBlue, anchor=east] at (-0.05,-0.15) {触摸};
\draw[dotpen=pcRed, wash=pcRed] (2.05,-0.15) circle (0.26);
\draw[pcRed, line width=0.7pt, -{Bar[width=6pt]}] (2.25,0.05) -- (3.2,0.66);
\node[lbl, text=pcRed, anchor=north] at (2.05,-0.45) {抑制};
\draw[dotpen=pcGray, wash=pcGray] (3.4,0.9) circle (0.34);
\node[lbl, anchor=south] at (3.4,1.3) {闸门};
\draw[pen] (3.75,0.9) -- (5.6,0.9); \draw[arr] (5.35,0.9) -- (5.65,0.9);
\node[lbl, anchor=west] at (5.7,0.9) {送往大脑};
}{脊髓里的闸门。疼痛信号经由闸门神经元进入大脑，而触摸唤醒一个抑制性神经元，把闸门掩上。所以揉一揉撞伤处有用，但只对中等程度的疼痛有效。}

\section*{来自上方的音量旋钮}

大脑自己就能把疼痛调大或调小。从上方的脑部通往闸门的，有几条广播频道——\nt{SERO}、\nt{NORA}，以及大脑自己分泌的、类似止痛药的特殊物质。在危险情境中，疼痛可能几乎消失：逃命的人没工夫瘸着走。对缓解的期待也会减轻疼痛，部分就是通过同样这些物质。

\section*{警报卡住的时候}

反复的疼痛会提高敏感度：同样的撞击开始变得更疼。这也是一种学习——一个在伤口愈合期间自己把自己拧紧的回路。但强回路有个危险：它可能“卡住”，在伤口早已愈合后还一直响个不停——就像第二章里那群自我维持的神经元。

\section*{老师与中断}

疼痛还是一位老师：对于多巴胺的“比预期的好还是差”，它起着负奖励的作用。它还会打断手头的事，就像一次去甲肾上腺素爆发：放下一切，先处理危险。而最快的反应——把手从烫的东西上缩回来——直接在脊髓里闭合，用的正是上一章那对对抗肌，你还没来得及感到疼，手就已经缩回来了。

\fullversion{14}
